\chapter{Conception et architecture du système}

\section{Architecture générale}

\noindent
\projectName{} adopte une architecture client-serveur composée de trois couches principales :

\begin{itemize}
    \item \textbf{Frontend React :} interface utilisateur, navigation, formulaires et affichage des données ;
    \item \textbf{Backend Laravel :} API REST, logique métier, authentification et validation ;
    \item \textbf{Base de données MySQL :} persistance des utilisateurs, offres, profils, candidatures et notifications.
\end{itemize}

\screenshotplaceholder{fig:architecture_smartjobs}{Schéma d'architecture globale de SmartJobs}

\section{Architecture frontend}

\noindent
Le frontend est organisé autour de pages et de composants réutilisables. Les pages publiques permettent de consulter les offres et d'accéder à l'authentification. Les espaces protégés sont séparés selon les rôles : candidat, recruteur et administrateur.

\noindent
Les composants partagés incluent notamment la barre de navigation, les cartes d'offres, les sélecteurs personnalisés, le panneau de notifications, les modales et le composant de discussion.

\section{Architecture backend}

\noindent
Le backend Laravel expose une API REST structurée. Les contrôleurs gèrent les actions principales :

\begin{itemize}
    \item authentification et profil utilisateur ;
    \item offres d'emploi ;
    \item candidatures et statuts ;
    \item quiz et questions ;
    \item notifications et messages ;
    \item paiement premium ;
    \item administration et modération.
\end{itemize}

\noindent
Les middlewares protègent les routes selon le rôle. Ainsi, un candidat ne peut pas accéder aux endpoints recruteur, et un recruteur ne peut pas accéder aux endpoints administrateur.

\section{Modèle de données}

\noindent
Le modèle de données repose sur plusieurs tables principales. Le tableau suivant résume les entités les plus importantes.

\begin{longtable}{|p{4cm}|p{9cm}|}
\hline
\textbf{Table} & \textbf{Rôle dans le système} \\
\hline
\endhead
users & Stocke les comptes, rôles, informations premium et quotas. \\
\hline
candidat\_profiles & Stocke les informations professionnelles du candidat, le CV et la photo. \\
\hline
recruteur\_profiles & Stocke les informations de l'établissement recruteur. \\
\hline
job\_offers & Stocke les offres publiées, leur statut, image et date d'expiration. \\
\hline
applications & Stocke les candidatures, statuts, scores de quiz et CV réutilisé. \\
\hline
quizzes / questions & Stockent les quiz QCM associés aux offres. \\
\hline
saved\_job\_offers & Stocke les offres sauvegardées par les candidats. \\
\hline
user\_notifications & Stocke les notifications internes. \\
\hline
application\_messages & Stocke les messages de discussion après acceptation. \\
\hline
\caption{Principales tables de la base de données}
\label{tab:tables_bd}
\end{longtable}

\screenshotplaceholder{fig:erd_smartjobs}{Modèle entité-association de la base de données SmartJobs}

\section{Flux de candidature}

\noindent
Le flux de candidature constitue une partie critique du système. Il doit garantir que le candidat possède un profil complet et un CV avant de postuler.

\begin{enumerate}
    \item Le candidat complète son profil et ajoute son CV.
    \item Il consulte une offre active.
    \item Il confirme sa candidature.
    \item Le backend crée une candidature unique pour le couple candidat/offre.
    \item Si l'offre contient un quiz, le candidat est dirigé vers la page du quiz.
    \item Le score du quiz met à jour la candidature existante.
\end{enumerate}

\screenshotplaceholder{fig:sequence_postulation}{Diagramme de séquence du flux de postulation}

\section{Sécurité et validation}

\noindent
Plusieurs mécanismes de sécurité ont été intégrés :

\begin{itemize}
    \item authentification par token avec Laravel Sanctum ;
    \item validation serveur des champs et des fichiers ;
    \item limitation des fichiers CV au format PDF ;
    \item limitation des images aux formats JPEG, PNG ou WEBP ;
    \item index unique pour éviter les candidatures dupliquées ;
    \item contrôle des rôles sur les routes sensibles.
\end{itemize}
